\documentclass[tikz,border=3pt]{standalone}
\input{tikz-tcee}
\begin{document}
\hyphenpenalty=10000 \exhyphenpenalty=10000
% Cómo se financia: los instrumentos de la cuenta financiera, de más estable a
% más volátil, y las 3 referencias de adecuación de las reservas (FMI).
\begin{tikzpicture}[tcee, font=\scriptsize,
  i/.style={draw=#1, fill=#1!14, rounded corners=2pt, line width=0.7pt, text width=1.45cm, align=center, minimum height=1.05cm, inner sep=2pt},
  r/.style={draw=dicze, fill=dicze!10, rounded corners=2pt, line width=0.7pt, text width=2.45cm, align=center, minimum height=1.2cm, inner sep=2pt}]
  % flecha de estabilidad (primero, debajo de las cajas)
  \fill[left color=dicfbcf!55, right color=dicparo!55] (-0.95,-0.85) rectangle (7.85,-0.65);
  \node[anchor=north west, font=\scriptsize\bfseries, text=dicfbcf] at (-0.95,-0.9) {más estable};
  \node[anchor=north east, font=\scriptsize\bfseries, text=dicparo] at (7.85,-0.9) {más volátil};
  % instrumentos
  \node[i=dicfbcf] at (0,0) {\textbf{Inversión directa}\\ trae tecnología};
  \node[i=dicexterna] at (1.8,0) {\textbf{Otra inversión}\\ préstamos, depósitos};
  \node[i=dicprestamos] at (3.6,0) {\textbf{Cartera}\\ efecto rebaño};
  \node[i=dicparo] at (5.4,0) {\textbf{Derivados}\\ aún más volátiles};
  \node[i=dicreferencia] at (7.2,0) {\textbf{Reservas}\\ ($\Delta R$)};
  % reservas
  \node[font=\footnotesize\bfseries, text=dicze, anchor=west] at (-0.95,-1.55) {¿Reservas suficientes? (orientativo, FMI)};
  \node[r] at (0.35,-2.45) {$\geq$ \textbf{3 meses} de importaciones\\ \emph{poco acceso a mercados}};
  \node[r] at (3.45,-2.45) {$\geq$ \textbf{100 \%} de la deuda a corto\\ \emph{\textsc{Greenspan}-\textsc{Guidotti}}};
  \node[r] at (6.55,-2.45) {$\geq$ \textbf{20 \%} de M2\\ \emph{fuga de depósitos}};
\end{tikzpicture}
\end{document}
